% =====================================================================
%  forecasting.tex  --  PHASE II Task 5 + PHASE V Task 12A
%
%  The forecasting comparison: seven models on one protocol, reported
%  as mean +/- sd over seeds, anchored on an untrained reference.
%
%  Intended placement: Section 5 (\label{sec:forecasting}), after the
%  state layer and before the decision layer.
%
%  Drop-in usage: % =====================================================================
%  forecasting.tex  --  PHASE II Task 5 + PHASE V Task 12A
%
%  The forecasting comparison: seven models on one protocol, reported
%  as mean +/- sd over seeds, anchored on an untrained reference.
%
%  Intended placement: Section 5 (\label{sec:forecasting}), after the
%  state layer and before the decision layer.
%
%  Drop-in usage: % =====================================================================
%  forecasting.tex  --  PHASE II Task 5 + PHASE V Task 12A
%
%  The forecasting comparison: seven models on one protocol, reported
%  as mean +/- sd over seeds, anchored on an untrained reference.
%
%  Intended placement: Section 5 (\label{sec:forecasting}), after the
%  state layer and before the decision layer.
%
%  Drop-in usage: % =====================================================================
%  forecasting.tex  --  PHASE II Task 5 + PHASE V Task 12A
%
%  The forecasting comparison: seven models on one protocol, reported
%  as mean +/- sd over seeds, anchored on an untrained reference.
%
%  Intended placement: Section 5 (\label{sec:forecasting}), after the
%  state layer and before the decision layer.
%
%  Drop-in usage: \input{forecasting}
%  Figures are referenced by bare name; add to the preamble
%      \graphicspath{{../outputs/phase2/figures/pelletizer-I/}
%                    {../outputs/phase6/figures/}}
%  Bibliography: paper/references.bib
%  Numbers below are generated by
%      python -m experiments.task5_forecasting_baselines   (single seed)
%      python run_phase5.py                                (seeds)
%  and stored in
%      outputs/phase2/task5/<machine>/comparison_table.csv
%      outputs/phase5/task12/<machine>/forecasting/seed_table.csv
%  Significance testing of every contrast quoted here is
%  Section~\ref{sec:significance}; attribution is
%  Section~\ref{sec:explain:window}.
% =====================================================================

\section{Forecasting the Non-Productive Horizon}
\label{sec:forecasting}

The state layer of Section~\ref{sec:states} produces a label per second.
A shutdown decision consumes something else: how much longer the machine
will remain non-productive. This section evaluates that forecasting step
as a supervised sequence task over the inferred labels, under one
protocol, against a reference that does no training at all.

\subsection{Protocol}
\label{sec:forecasting:protocol}

The labelled record is decimated $10\times$ to $0.1$\,Hz and cut into
windows of $600$\,s lookback and $300$\,s horizon at a stride of
$120$\,s. The split is chronological ($70/15/15$) and windows are built
\emph{inside} each split, so no window straddles a split boundary and no
future sample can enter a training window. Class weights are balanced
identically for every model. The test set is $5\,449$ windows, on which
the mean true STANDBY content is $20.1$\,s of a $300$\,s horizon and
only $13.5\%$ of windows contain any STANDBY at all.

Two metrics are reported. STANDBY $F_1$ is the per-step classification
quality on the minority class the decision layer cares about. The
per-window \emph{STANDBY-seconds MAE} --- the absolute error in how many
of the next $300$\,s are predicted to be standby --- is the quantity
closer to what a duration-consuming policy would use, and it is the
metric on which the significance tests of
Section~\ref{sec:significance} are pre-registered.

\paragraph{Seeds are reported, not a single run.}
Every stochastic model is trained at multiple seeds and the table
reports mean\,$\pm$\,sd: $15$ seeds for the gradient-boosted tree, $5$
for each neural architecture, $1$ for the deterministic reference. This
is not a formality here. The spread between a good and a bad seed of one
architecture reaches $0.166$ of $F_1$, which is larger than the entire
spread between architectures, so a single-seed table of this comparison
is not a ranking (Figure~\ref{fig:seed-stability}).

\paragraph{The reference is not a trained model.}
\textsc{persistence} continues the last observed state across the whole
horizon. It has no parameters and no fitting step. A comparison among
trained models can establish which architecture is better; only a
comparison against not training can establish that the forecasting stage
belongs in the framework at all. This row was absent from the
comparison as first run, and adding it changed the conclusion.

\subsection{Results}
\label{sec:forecasting:results}

\begin{table}[t]
\centering
\caption{Forecasting comparison on $5\,449$ held-out windows of
pelletizer~I, mean\,$\pm$\,sd over seeds. Lower MAE is better. The
untrained reference is deterministic. Ordered by MAE.}
\label{tab:forecasting}
\small
\setlength{\tabcolsep}{4pt}
\begin{tabular}{@{}lrrrr@{}}
\toprule
Model & seeds & STANDBY $F_1$ & MAE (s) & params \\
\midrule
XGBoost \cite{chen2016xgboost}          & 15 & $\mathbf{0.689 \pm 0.003}$ & $\mathbf{11.50 \pm 0.11}$ & 236 trees \\
\textsc{persistence} (untrained)        &  1 & $0.681$                    & $12.95$                   & $\mathbf{0}$ \\
Transformer \cite{nie2023patchtst}      &  5 & $0.595 \pm 0.032$          & $17.45 \pm 3.13$          & $310\,488$ \\
GRU \cite{cho2014gru}                   &  5 & $0.599 \pm 0.073$          & $17.68 \pm 5.59$          & $121\,604$ \\
TCN \cite{bai2018tcn}                   &  5 & $0.588 \pm 0.050$          & $18.66 \pm 3.29$          & $131\,332$ \\
Vanilla LSTM \cite{hochreiter1997lstm}  &  5 & $0.578 \pm 0.036$          & $19.12 \pm 3.01$          & $93\,816$ \\
Seq2Seq LSTM \cite{sutskever2014seq2seq}&  5 & $0.561 \pm 0.064$          & $20.90 \pm 4.74$          & $161\,028$ \\
\bottomrule
\end{tabular}
\end{table}

Table~\ref{tab:forecasting} orders the seven models by the error metric
the decision layer would consume. Three results follow, and the first
two are not the ones this comparison was designed to produce.

\textbf{Every neural architecture loses to the untrained reference.}
Persistence reaches $12.95$\,s of per-window MAE against $17.45$--$20.90$\,s
for the five sequence models, and $F_1 = 0.681$ against
$0.561$--$0.599$. The margins are $4.5$--$8.0$\,s of MAE, all of them
above the $1$\,s threshold fixed in advance as the smallest difference
that could change a decision, and all significant at
$p_{\text{Holm}}\le6\times10^{-31}$ (Section~\ref{sec:significance}). The
sequence-to-sequence LSTM --- the architecture this framework originally
proposed for this stage --- is the worst of the seven, losing to
persistence by $7.95$\,s of MAE while costing $161$k parameters and
$865\pm75$\,s of training.

\textbf{One trained model clears the reference, and it is the tree.}
XGBoost reaches $11.50$\,s, $1.45$\,s better than persistence at
$p_{\text{Holm}}=7\times10^{-4}$: statistically detectable, above the
practical threshold, and the only comparison in the family where
training beats not training. It is also the cheapest of the trained
models at inference ($0.18$\,ms per window against $0.51$--$1.16$\,ms
for the neural models). \textbf{The gradient-boosted tree is therefore the
forecasting component of the proposed framework}, and the sequence
models are reported as what they are: compared baselines that did not
earn the stage.

\textbf{The tree is also an order of magnitude more stable.} Its $F_1$
standard deviation over $15$ seeds is $0.003$; the neural models' over
$5$ seeds is $0.032$--$0.073$, a factor of $13$--$29$
(Figure~\ref{fig:seed-stability}). The practical consequence is visible
in the same figure: seed~$0$ is the best of five draws for four of the
five neural architectures, so a single-seed protocol --- which is what
the first run of this comparison used --- reports each of them at close
to its most favourable outcome. That protocol was sound and its runs
reproduce bit-for-bit; it simply cannot support a ranking.

\begin{figure}[t]
  \centering
  \includegraphics[width=\linewidth]{fig_p6_seed_stability}
  \caption{STANDBY $F_1$ by training seed for all seven forecasters.
  The tree's $15$ seeds span $0.009$ of $F_1$; the neural
  architectures' five seeds span $0.081$--$0.166$, wider than the gap
  between architectures. The dashed line is the untrained reference.}
  \label{fig:seed-stability}
\end{figure}

\subsection{Why the untrained reference is so hard to beat}
\label{sec:forecasting:why}

The result is not an artefact of an unfair protocol, and the explanation
is available from two independent directions.

The first is the operating regime. The horizon is $300$\,s while the
median state dwell after the minimum-dwell constraint is $92.5$\,s with
a heavy tail --- STANDBY runs last minutes to hours
(Section~\ref{sec:results:task4}). ``Nothing changes in the next five
minutes'' is therefore right most of the time, and a forecaster deployed
in this regime has to beat that, not a strawman.

The second is measured on the winning model itself. Attributing
XGBoost's predictions (Section~\ref{sec:explain:window}) puts $43.5\%$
of the attribution on the final observed value of the lookback window,
against $26.4\%$ on block means and $26.3\%$ on block standard
deviations. A model whose largest single input is the state the machine
is in \emph{now} is a refinement of persistence, and cannot be expected
to beat it by much. It does not: $1.45$\,s. The margin is real, and it
is small for a reason the attribution makes explicit rather than leaving
to speculation.

\paragraph{A caveat for deployment elsewhere.}
The same attribution shows that supply voltage is the largest single
channel contributor to this model, at $20.0\%$ across its nine design
columns and $11.8\%$ on the last observed value alone. Supply voltage is
not a state variable of the machine: it sags under load and recovers
when the drive unloads, so it is a real but indirect indicator, and one
tied to this site's point of common coupling and transformer impedance.
\textbf{A deployment on another site should not expect that contribution
to transfer}, and the cross-dataset failure reported in
Section~\ref{sec:crossmachine:dataset} is consistent with it. Split gain
ranks voltage tenth of the twenty-five channels and gives no such
warning, which is one more reason the attribution is reported.

\paragraph{Scope.}
This comparison is one machine, one horizon, and one decimation.
Section~\ref{sec:crossmachine} repeats the training across all eight
machines and asks whether any of it transfers; the answer bears directly
on how much of Table~\ref{tab:forecasting} is a property of the method
and how much of pelletizer~I.

%  Figures are referenced by bare name; add to the preamble
%      \graphicspath{{../outputs/phase2/figures/pelletizer-I/}
%                    {../outputs/phase6/figures/}}
%  Bibliography: paper/references.bib
%  Numbers below are generated by
%      python -m experiments.task5_forecasting_baselines   (single seed)
%      python run_phase5.py                                (seeds)
%  and stored in
%      outputs/phase2/task5/<machine>/comparison_table.csv
%      outputs/phase5/task12/<machine>/forecasting/seed_table.csv
%  Significance testing of every contrast quoted here is
%  Section~\ref{sec:significance}; attribution is
%  Section~\ref{sec:explain:window}.
% =====================================================================

\section{Forecasting the Non-Productive Horizon}
\label{sec:forecasting}

The state layer of Section~\ref{sec:states} produces a label per second.
A shutdown decision consumes something else: how much longer the machine
will remain non-productive. This section evaluates that forecasting step
as a supervised sequence task over the inferred labels, under one
protocol, against a reference that does no training at all.

\subsection{Protocol}
\label{sec:forecasting:protocol}

The labelled record is decimated $10\times$ to $0.1$\,Hz and cut into
windows of $600$\,s lookback and $300$\,s horizon at a stride of
$120$\,s. The split is chronological ($70/15/15$) and windows are built
\emph{inside} each split, so no window straddles a split boundary and no
future sample can enter a training window. Class weights are balanced
identically for every model. The test set is $5\,449$ windows, on which
the mean true STANDBY content is $20.1$\,s of a $300$\,s horizon and
only $13.5\%$ of windows contain any STANDBY at all.

Two metrics are reported. STANDBY $F_1$ is the per-step classification
quality on the minority class the decision layer cares about. The
per-window \emph{STANDBY-seconds MAE} --- the absolute error in how many
of the next $300$\,s are predicted to be standby --- is the quantity
closer to what a duration-consuming policy would use, and it is the
metric on which the significance tests of
Section~\ref{sec:significance} are pre-registered.

\paragraph{Seeds are reported, not a single run.}
Every stochastic model is trained at multiple seeds and the table
reports mean\,$\pm$\,sd: $15$ seeds for the gradient-boosted tree, $5$
for each neural architecture, $1$ for the deterministic reference. This
is not a formality here. The spread between a good and a bad seed of one
architecture reaches $0.166$ of $F_1$, which is larger than the entire
spread between architectures, so a single-seed table of this comparison
is not a ranking (Figure~\ref{fig:seed-stability}).

\paragraph{The reference is not a trained model.}
\textsc{persistence} continues the last observed state across the whole
horizon. It has no parameters and no fitting step. A comparison among
trained models can establish which architecture is better; only a
comparison against not training can establish that the forecasting stage
belongs in the framework at all. This row was absent from the
comparison as first run, and adding it changed the conclusion.

\subsection{Results}
\label{sec:forecasting:results}

\begin{table}[t]
\centering
\caption{Forecasting comparison on $5\,449$ held-out windows of
pelletizer~I, mean\,$\pm$\,sd over seeds. Lower MAE is better. The
untrained reference is deterministic. Ordered by MAE.}
\label{tab:forecasting}
\small
\setlength{\tabcolsep}{4pt}
\begin{tabular}{@{}lrrrr@{}}
\toprule
Model & seeds & STANDBY $F_1$ & MAE (s) & params \\
\midrule
XGBoost \cite{chen2016xgboost}          & 15 & $\mathbf{0.689 \pm 0.003}$ & $\mathbf{11.50 \pm 0.11}$ & 236 trees \\
\textsc{persistence} (untrained)        &  1 & $0.681$                    & $12.95$                   & $\mathbf{0}$ \\
Transformer \cite{nie2023patchtst}      &  5 & $0.595 \pm 0.032$          & $17.45 \pm 3.13$          & $310\,488$ \\
GRU \cite{cho2014gru}                   &  5 & $0.599 \pm 0.073$          & $17.68 \pm 5.59$          & $121\,604$ \\
TCN \cite{bai2018tcn}                   &  5 & $0.588 \pm 0.050$          & $18.66 \pm 3.29$          & $131\,332$ \\
Vanilla LSTM \cite{hochreiter1997lstm}  &  5 & $0.578 \pm 0.036$          & $19.12 \pm 3.01$          & $93\,816$ \\
Seq2Seq LSTM \cite{sutskever2014seq2seq}&  5 & $0.561 \pm 0.064$          & $20.90 \pm 4.74$          & $161\,028$ \\
\bottomrule
\end{tabular}
\end{table}

Table~\ref{tab:forecasting} orders the seven models by the error metric
the decision layer would consume. Three results follow, and the first
two are not the ones this comparison was designed to produce.

\textbf{Every neural architecture loses to the untrained reference.}
Persistence reaches $12.95$\,s of per-window MAE against $17.45$--$20.90$\,s
for the five sequence models, and $F_1 = 0.681$ against
$0.561$--$0.599$. The margins are $4.5$--$8.0$\,s of MAE, all of them
above the $1$\,s threshold fixed in advance as the smallest difference
that could change a decision, and all significant at
$p_{\text{Holm}}\le6\times10^{-31}$ (Section~\ref{sec:significance}). The
sequence-to-sequence LSTM --- the architecture this framework originally
proposed for this stage --- is the worst of the seven, losing to
persistence by $7.95$\,s of MAE while costing $161$k parameters and
$865\pm75$\,s of training.

\textbf{One trained model clears the reference, and it is the tree.}
XGBoost reaches $11.50$\,s, $1.45$\,s better than persistence at
$p_{\text{Holm}}=7\times10^{-4}$: statistically detectable, above the
practical threshold, and the only comparison in the family where
training beats not training. It is also the cheapest of the trained
models at inference ($0.18$\,ms per window against $0.51$--$1.16$\,ms
for the neural models). \textbf{The gradient-boosted tree is therefore the
forecasting component of the proposed framework}, and the sequence
models are reported as what they are: compared baselines that did not
earn the stage.

\textbf{The tree is also an order of magnitude more stable.} Its $F_1$
standard deviation over $15$ seeds is $0.003$; the neural models' over
$5$ seeds is $0.032$--$0.073$, a factor of $13$--$29$
(Figure~\ref{fig:seed-stability}). The practical consequence is visible
in the same figure: seed~$0$ is the best of five draws for four of the
five neural architectures, so a single-seed protocol --- which is what
the first run of this comparison used --- reports each of them at close
to its most favourable outcome. That protocol was sound and its runs
reproduce bit-for-bit; it simply cannot support a ranking.

\begin{figure}[t]
  \centering
  \includegraphics[width=\linewidth]{fig_p6_seed_stability}
  \caption{STANDBY $F_1$ by training seed for all seven forecasters.
  The tree's $15$ seeds span $0.009$ of $F_1$; the neural
  architectures' five seeds span $0.081$--$0.166$, wider than the gap
  between architectures. The dashed line is the untrained reference.}
  \label{fig:seed-stability}
\end{figure}

\subsection{Why the untrained reference is so hard to beat}
\label{sec:forecasting:why}

The result is not an artefact of an unfair protocol, and the explanation
is available from two independent directions.

The first is the operating regime. The horizon is $300$\,s while the
median state dwell after the minimum-dwell constraint is $92.5$\,s with
a heavy tail --- STANDBY runs last minutes to hours
(Section~\ref{sec:results:task4}). ``Nothing changes in the next five
minutes'' is therefore right most of the time, and a forecaster deployed
in this regime has to beat that, not a strawman.

The second is measured on the winning model itself. Attributing
XGBoost's predictions (Section~\ref{sec:explain:window}) puts $43.5\%$
of the attribution on the final observed value of the lookback window,
against $26.4\%$ on block means and $26.3\%$ on block standard
deviations. A model whose largest single input is the state the machine
is in \emph{now} is a refinement of persistence, and cannot be expected
to beat it by much. It does not: $1.45$\,s. The margin is real, and it
is small for a reason the attribution makes explicit rather than leaving
to speculation.

\paragraph{A caveat for deployment elsewhere.}
The same attribution shows that supply voltage is the largest single
channel contributor to this model, at $20.0\%$ across its nine design
columns and $11.8\%$ on the last observed value alone. Supply voltage is
not a state variable of the machine: it sags under load and recovers
when the drive unloads, so it is a real but indirect indicator, and one
tied to this site's point of common coupling and transformer impedance.
\textbf{A deployment on another site should not expect that contribution
to transfer}, and the cross-dataset failure reported in
Section~\ref{sec:crossmachine:dataset} is consistent with it. Split gain
ranks voltage tenth of the twenty-five channels and gives no such
warning, which is one more reason the attribution is reported.

\paragraph{Scope.}
This comparison is one machine, one horizon, and one decimation.
Section~\ref{sec:crossmachine} repeats the training across all eight
machines and asks whether any of it transfers; the answer bears directly
on how much of Table~\ref{tab:forecasting} is a property of the method
and how much of pelletizer~I.

%  Figures are referenced by bare name; add to the preamble
%      \graphicspath{{../outputs/phase2/figures/pelletizer-I/}
%                    {../outputs/phase6/figures/}}
%  Bibliography: paper/references.bib
%  Numbers below are generated by
%      python -m experiments.task5_forecasting_baselines   (single seed)
%      python run_phase5.py                                (seeds)
%  and stored in
%      outputs/phase2/task5/<machine>/comparison_table.csv
%      outputs/phase5/task12/<machine>/forecasting/seed_table.csv
%  Significance testing of every contrast quoted here is
%  Section~\ref{sec:significance}; attribution is
%  Section~\ref{sec:explain:window}.
% =====================================================================

\section{Forecasting the Non-Productive Horizon}
\label{sec:forecasting}

The state layer of Section~\ref{sec:states} produces a label per second.
A shutdown decision consumes something else: how much longer the machine
will remain non-productive. This section evaluates that forecasting step
as a supervised sequence task over the inferred labels, under one
protocol, against a reference that does no training at all.

\subsection{Protocol}
\label{sec:forecasting:protocol}

The labelled record is decimated $10\times$ to $0.1$\,Hz and cut into
windows of $600$\,s lookback and $300$\,s horizon at a stride of
$120$\,s. The split is chronological ($70/15/15$) and windows are built
\emph{inside} each split, so no window straddles a split boundary and no
future sample can enter a training window. Class weights are balanced
identically for every model. The test set is $5\,449$ windows, on which
the mean true STANDBY content is $20.1$\,s of a $300$\,s horizon and
only $13.5\%$ of windows contain any STANDBY at all.

Two metrics are reported. STANDBY $F_1$ is the per-step classification
quality on the minority class the decision layer cares about. The
per-window \emph{STANDBY-seconds MAE} --- the absolute error in how many
of the next $300$\,s are predicted to be standby --- is the quantity
closer to what a duration-consuming policy would use, and it is the
metric on which the significance tests of
Section~\ref{sec:significance} are pre-registered.

\paragraph{Seeds are reported, not a single run.}
Every stochastic model is trained at multiple seeds and the table
reports mean\,$\pm$\,sd: $15$ seeds for the gradient-boosted tree, $5$
for each neural architecture, $1$ for the deterministic reference. This
is not a formality here. The spread between a good and a bad seed of one
architecture reaches $0.166$ of $F_1$, which is larger than the entire
spread between architectures, so a single-seed table of this comparison
is not a ranking (Figure~\ref{fig:seed-stability}).

\paragraph{The reference is not a trained model.}
\textsc{persistence} continues the last observed state across the whole
horizon. It has no parameters and no fitting step. A comparison among
trained models can establish which architecture is better; only a
comparison against not training can establish that the forecasting stage
belongs in the framework at all. This row was absent from the
comparison as first run, and adding it changed the conclusion.

\subsection{Results}
\label{sec:forecasting:results}

\begin{table}[t]
\centering
\caption{Forecasting comparison on $5\,449$ held-out windows of
pelletizer~I, mean\,$\pm$\,sd over seeds. Lower MAE is better. The
untrained reference is deterministic. Ordered by MAE.}
\label{tab:forecasting}
\small
\setlength{\tabcolsep}{4pt}
\begin{tabular}{@{}lrrrr@{}}
\toprule
Model & seeds & STANDBY $F_1$ & MAE (s) & params \\
\midrule
XGBoost \cite{chen2016xgboost}          & 15 & $\mathbf{0.689 \pm 0.003}$ & $\mathbf{11.50 \pm 0.11}$ & 236 trees \\
\textsc{persistence} (untrained)        &  1 & $0.681$                    & $12.95$                   & $\mathbf{0}$ \\
Transformer \cite{nie2023patchtst}      &  5 & $0.595 \pm 0.032$          & $17.45 \pm 3.13$          & $310\,488$ \\
GRU \cite{cho2014gru}                   &  5 & $0.599 \pm 0.073$          & $17.68 \pm 5.59$          & $121\,604$ \\
TCN \cite{bai2018tcn}                   &  5 & $0.588 \pm 0.050$          & $18.66 \pm 3.29$          & $131\,332$ \\
Vanilla LSTM \cite{hochreiter1997lstm}  &  5 & $0.578 \pm 0.036$          & $19.12 \pm 3.01$          & $93\,816$ \\
Seq2Seq LSTM \cite{sutskever2014seq2seq}&  5 & $0.561 \pm 0.064$          & $20.90 \pm 4.74$          & $161\,028$ \\
\bottomrule
\end{tabular}
\end{table}

Table~\ref{tab:forecasting} orders the seven models by the error metric
the decision layer would consume. Three results follow, and the first
two are not the ones this comparison was designed to produce.

\textbf{Every neural architecture loses to the untrained reference.}
Persistence reaches $12.95$\,s of per-window MAE against $17.45$--$20.90$\,s
for the five sequence models, and $F_1 = 0.681$ against
$0.561$--$0.599$. The margins are $4.5$--$8.0$\,s of MAE, all of them
above the $1$\,s threshold fixed in advance as the smallest difference
that could change a decision, and all significant at
$p_{\text{Holm}}\le6\times10^{-31}$ (Section~\ref{sec:significance}). The
sequence-to-sequence LSTM --- the architecture this framework originally
proposed for this stage --- is the worst of the seven, losing to
persistence by $7.95$\,s of MAE while costing $161$k parameters and
$865\pm75$\,s of training.

\textbf{One trained model clears the reference, and it is the tree.}
XGBoost reaches $11.50$\,s, $1.45$\,s better than persistence at
$p_{\text{Holm}}=7\times10^{-4}$: statistically detectable, above the
practical threshold, and the only comparison in the family where
training beats not training. It is also the cheapest of the trained
models at inference ($0.18$\,ms per window against $0.51$--$1.16$\,ms
for the neural models). \textbf{The gradient-boosted tree is therefore the
forecasting component of the proposed framework}, and the sequence
models are reported as what they are: compared baselines that did not
earn the stage.

\textbf{The tree is also an order of magnitude more stable.} Its $F_1$
standard deviation over $15$ seeds is $0.003$; the neural models' over
$5$ seeds is $0.032$--$0.073$, a factor of $13$--$29$
(Figure~\ref{fig:seed-stability}). The practical consequence is visible
in the same figure: seed~$0$ is the best of five draws for four of the
five neural architectures, so a single-seed protocol --- which is what
the first run of this comparison used --- reports each of them at close
to its most favourable outcome. That protocol was sound and its runs
reproduce bit-for-bit; it simply cannot support a ranking.

\begin{figure}[t]
  \centering
  \includegraphics[width=\linewidth]{fig_p6_seed_stability}
  \caption{STANDBY $F_1$ by training seed for all seven forecasters.
  The tree's $15$ seeds span $0.009$ of $F_1$; the neural
  architectures' five seeds span $0.081$--$0.166$, wider than the gap
  between architectures. The dashed line is the untrained reference.}
  \label{fig:seed-stability}
\end{figure}

\subsection{Why the untrained reference is so hard to beat}
\label{sec:forecasting:why}

The result is not an artefact of an unfair protocol, and the explanation
is available from two independent directions.

The first is the operating regime. The horizon is $300$\,s while the
median state dwell after the minimum-dwell constraint is $92.5$\,s with
a heavy tail --- STANDBY runs last minutes to hours
(Section~\ref{sec:results:task4}). ``Nothing changes in the next five
minutes'' is therefore right most of the time, and a forecaster deployed
in this regime has to beat that, not a strawman.

The second is measured on the winning model itself. Attributing
XGBoost's predictions (Section~\ref{sec:explain:window}) puts $43.5\%$
of the attribution on the final observed value of the lookback window,
against $26.4\%$ on block means and $26.3\%$ on block standard
deviations. A model whose largest single input is the state the machine
is in \emph{now} is a refinement of persistence, and cannot be expected
to beat it by much. It does not: $1.45$\,s. The margin is real, and it
is small for a reason the attribution makes explicit rather than leaving
to speculation.

\paragraph{A caveat for deployment elsewhere.}
The same attribution shows that supply voltage is the largest single
channel contributor to this model, at $20.0\%$ across its nine design
columns and $11.8\%$ on the last observed value alone. Supply voltage is
not a state variable of the machine: it sags under load and recovers
when the drive unloads, so it is a real but indirect indicator, and one
tied to this site's point of common coupling and transformer impedance.
\textbf{A deployment on another site should not expect that contribution
to transfer}, and the cross-dataset failure reported in
Section~\ref{sec:crossmachine:dataset} is consistent with it. Split gain
ranks voltage tenth of the twenty-five channels and gives no such
warning, which is one more reason the attribution is reported.

\paragraph{Scope.}
This comparison is one machine, one horizon, and one decimation.
Section~\ref{sec:crossmachine} repeats the training across all eight
machines and asks whether any of it transfers; the answer bears directly
on how much of Table~\ref{tab:forecasting} is a property of the method
and how much of pelletizer~I.

%  Figures are referenced by bare name; add to the preamble
%      \graphicspath{{../outputs/phase2/figures/pelletizer-I/}
%                    {../outputs/phase6/figures/}}
%  Bibliography: paper/references.bib
%  Numbers below are generated by
%      python -m experiments.task5_forecasting_baselines   (single seed)
%      python run_phase5.py                                (seeds)
%  and stored in
%      outputs/phase2/task5/<machine>/comparison_table.csv
%      outputs/phase5/task12/<machine>/forecasting/seed_table.csv
%  Significance testing of every contrast quoted here is
%  Section~\ref{sec:significance}; attribution is
%  Section~\ref{sec:explain:window}.
% =====================================================================

\section{Forecasting the Non-Productive Horizon}
\label{sec:forecasting}

The state layer of Section~\ref{sec:states} produces a label per second.
A shutdown decision consumes something else: how much longer the machine
will remain non-productive. This section evaluates that forecasting step
as a supervised sequence task over the inferred labels, under one
protocol, against a reference that does no training at all.

\subsection{Protocol}
\label{sec:forecasting:protocol}

The labelled record is decimated $10\times$ to $0.1$\,Hz and cut into
windows of $600$\,s lookback and $300$\,s horizon at a stride of
$120$\,s. The split is chronological ($70/15/15$) and windows are built
\emph{inside} each split, so no window straddles a split boundary and no
future sample can enter a training window. Class weights are balanced
identically for every model. The test set is $5\,449$ windows, on which
the mean true STANDBY content is $20.1$\,s of a $300$\,s horizon and
only $13.5\%$ of windows contain any STANDBY at all.

Two metrics are reported. STANDBY $F_1$ is the per-step classification
quality on the minority class the decision layer cares about. The
per-window \emph{STANDBY-seconds MAE} --- the absolute error in how many
of the next $300$\,s are predicted to be standby --- is the quantity
closer to what a duration-consuming policy would use, and it is the
metric on which the significance tests of
Section~\ref{sec:significance} are pre-registered.

\paragraph{Seeds are reported, not a single run.}
Every stochastic model is trained at multiple seeds and the table
reports mean\,$\pm$\,sd: $15$ seeds for the gradient-boosted tree, $5$
for each neural architecture, $1$ for the deterministic reference. This
is not a formality here. The spread between a good and a bad seed of one
architecture reaches $0.166$ of $F_1$, which is larger than the entire
spread between architectures, so a single-seed table of this comparison
is not a ranking (Figure~\ref{fig:seed-stability}).

\paragraph{The reference is not a trained model.}
\textsc{persistence} continues the last observed state across the whole
horizon. It has no parameters and no fitting step. A comparison among
trained models can establish which architecture is better; only a
comparison against not training can establish that the forecasting stage
belongs in the framework at all. This row was absent from the
comparison as first run, and adding it changed the conclusion.

\subsection{Results}
\label{sec:forecasting:results}

\begin{table}[t]
\centering
\caption{Forecasting comparison on $5\,449$ held-out windows of
pelletizer~I, mean\,$\pm$\,sd over seeds. Lower MAE is better. The
untrained reference is deterministic. Ordered by MAE.}
\label{tab:forecasting}
\small
\setlength{\tabcolsep}{4pt}
\begin{tabular}{@{}lrrrr@{}}
\toprule
Model & seeds & STANDBY $F_1$ & MAE (s) & params \\
\midrule
XGBoost \cite{chen2016xgboost}          & 15 & $\mathbf{0.689 \pm 0.003}$ & $\mathbf{11.50 \pm 0.11}$ & 236 trees \\
\textsc{persistence} (untrained)        &  1 & $0.681$                    & $12.95$                   & $\mathbf{0}$ \\
Transformer \cite{nie2023patchtst}      &  5 & $0.595 \pm 0.032$          & $17.45 \pm 3.13$          & $310\,488$ \\
GRU \cite{cho2014gru}                   &  5 & $0.599 \pm 0.073$          & $17.68 \pm 5.59$          & $121\,604$ \\
TCN \cite{bai2018tcn}                   &  5 & $0.588 \pm 0.050$          & $18.66 \pm 3.29$          & $131\,332$ \\
Vanilla LSTM \cite{hochreiter1997lstm}  &  5 & $0.578 \pm 0.036$          & $19.12 \pm 3.01$          & $93\,816$ \\
Seq2Seq LSTM \cite{sutskever2014seq2seq}&  5 & $0.561 \pm 0.064$          & $20.90 \pm 4.74$          & $161\,028$ \\
\bottomrule
\end{tabular}
\end{table}

Table~\ref{tab:forecasting} orders the seven models by the error metric
the decision layer would consume. Three results follow, and the first
two are not the ones this comparison was designed to produce.

\textbf{Every neural architecture loses to the untrained reference.}
Persistence reaches $12.95$\,s of per-window MAE against $17.45$--$20.90$\,s
for the five sequence models, and $F_1 = 0.681$ against
$0.561$--$0.599$. The margins are $4.5$--$8.0$\,s of MAE, all of them
above the $1$\,s threshold fixed in advance as the smallest difference
that could change a decision, and all significant at
$p_{\text{Holm}}\le6\times10^{-31}$ (Section~\ref{sec:significance}). The
sequence-to-sequence LSTM --- the architecture this framework originally
proposed for this stage --- is the worst of the seven, losing to
persistence by $7.95$\,s of MAE while costing $161$k parameters and
$865\pm75$\,s of training.

\textbf{One trained model clears the reference, and it is the tree.}
XGBoost reaches $11.50$\,s, $1.45$\,s better than persistence at
$p_{\text{Holm}}=7\times10^{-4}$: statistically detectable, above the
practical threshold, and the only comparison in the family where
training beats not training. It is also the cheapest of the trained
models at inference ($0.18$\,ms per window against $0.51$--$1.16$\,ms
for the neural models). \textbf{The gradient-boosted tree is therefore the
forecasting component of the proposed framework}, and the sequence
models are reported as what they are: compared baselines that did not
earn the stage.

\textbf{The tree is also an order of magnitude more stable.} Its $F_1$
standard deviation over $15$ seeds is $0.003$; the neural models' over
$5$ seeds is $0.032$--$0.073$, a factor of $13$--$29$
(Figure~\ref{fig:seed-stability}). The practical consequence is visible
in the same figure: seed~$0$ is the best of five draws for four of the
five neural architectures, so a single-seed protocol --- which is what
the first run of this comparison used --- reports each of them at close
to its most favourable outcome. That protocol was sound and its runs
reproduce bit-for-bit; it simply cannot support a ranking.

\begin{figure}[t]
  \centering
  \includegraphics[width=\linewidth]{fig_p6_seed_stability}
  \caption{STANDBY $F_1$ by training seed for all seven forecasters.
  The tree's $15$ seeds span $0.009$ of $F_1$; the neural
  architectures' five seeds span $0.081$--$0.166$, wider than the gap
  between architectures. The dashed line is the untrained reference.}
  \label{fig:seed-stability}
\end{figure}

\subsection{Why the untrained reference is so hard to beat}
\label{sec:forecasting:why}

The result is not an artefact of an unfair protocol, and the explanation
is available from two independent directions.

The first is the operating regime. The horizon is $300$\,s while the
median state dwell after the minimum-dwell constraint is $92.5$\,s with
a heavy tail --- STANDBY runs last minutes to hours
(Section~\ref{sec:results:task4}). ``Nothing changes in the next five
minutes'' is therefore right most of the time, and a forecaster deployed
in this regime has to beat that, not a strawman.

The second is measured on the winning model itself. Attributing
XGBoost's predictions (Section~\ref{sec:explain:window}) puts $43.5\%$
of the attribution on the final observed value of the lookback window,
against $26.4\%$ on block means and $26.3\%$ on block standard
deviations. A model whose largest single input is the state the machine
is in \emph{now} is a refinement of persistence, and cannot be expected
to beat it by much. It does not: $1.45$\,s. The margin is real, and it
is small for a reason the attribution makes explicit rather than leaving
to speculation.

\paragraph{A caveat for deployment elsewhere.}
The same attribution shows that supply voltage is the largest single
channel contributor to this model, at $20.0\%$ across its nine design
columns and $11.8\%$ on the last observed value alone. Supply voltage is
not a state variable of the machine: it sags under load and recovers
when the drive unloads, so it is a real but indirect indicator, and one
tied to this site's point of common coupling and transformer impedance.
\textbf{A deployment on another site should not expect that contribution
to transfer}, and the cross-dataset failure reported in
Section~\ref{sec:crossmachine:dataset} is consistent with it. Split gain
ranks voltage tenth of the twenty-five channels and gives no such
warning, which is one more reason the attribution is reported.

\paragraph{Scope.}
This comparison is one machine, one horizon, and one decimation.
Section~\ref{sec:crossmachine} repeats the training across all eight
machines and asks whether any of it transfers; the answer bears directly
on how much of Table~\ref{tab:forecasting} is a property of the method
and how much of pelletizer~I.
